\begin{summarybox}
	\textbf{\textcolor{CSummary}{一句话核心}}
	\textbf{椭圆参数方程将坐标用三角函数表示，便于计算最值和轨迹。}

	\medskip
	\textbf{\textcolor{CSummary}{使用条件}}
	\begin{itemize}[leftmargin=*, itemsep=0.5em, topsep=0.4em]
		\item \textbf{条件 1（标准形式）：} 椭圆方程为 $\frac{x^2}{a^2}+\frac{y^2}{b^2}=1\;(a>b>0)$，中心在原点，焦点在 $x$ 轴。
		\item \textbf{条件 2（参数范围）：} 参数 $\theta$ 通常取 $[0,2\pi)$，实际问题中根据要求选定。
	\end{itemize}

	\medskip
	\textbf{\textcolor{CSummary}{距离最值的两种观察方向}}
	\begin{itemize}[leftmargin=*, itemsep=0.5em, topsep=0.4em]
		\item \textbf{由内到外（固定点 $O$ $\to$ 椭圆）：} 固定圆心 $O$，使距离圆 $x^2+y^2=r^2$ 的半径逐渐增大；在 $r=b$ 时首次与椭圆相切，在 $r=a$ 时最后一次相切。
		\item \textbf{由外到内（椭圆动点 $P$ $\to$ 固定点 $O$）：} 设 $P(a\cos\theta,b\sin\theta)$，计算 $OP^2=b^2+(a^2-b^2)\cos^2\theta$，由三角函数有界性求最值。
	\end{itemize}
	\textbf{共同结论：} $\boxed{OP_{\min}=b,\quad OP_{\max}=a}$；一个从距离圆找临界相切，一个从动点求距离函数最值。

	\medskip
	\textbf{\textcolor{CSummary}{关键提醒}}
	\begin{itemize}[leftmargin=*, itemsep=0.5em, topsep=0.4em]
		\item \textbf{易错点（离心角非极角）：} 离心角 $\theta$ 与极角 $\varphi$ 不同，关系为 $\tan\varphi=\frac{b}{a}\tan\theta$。
		\item \textbf{检查项（有界性验证）：} 使用 $\sin\theta,\cos\theta$ 有界性求最值时，需验证取等条件。
	\end{itemize}

\end{summarybox}
